\section{Fully Homomorphic Encryption}
\label{sec:fhe}
\gls{he} allows computation directly on encrypted data: an untrusted party can evaluate a function over ciphertexts without learning anything about the underlying plaintexts. % Schemes are distinguished by the circuits they can evaluate --- \emph{leveled} schemes support circuits up to a fixed multiplicative depth, whereas true \gls{fhe} supports arbitrary depth. 
With few exceptions, the best-performing schemes \cite{cryptoeprint:2011/277, cryptoeprint:2012/144, cryptoeprint:2016/421, cryptoeprint:2018/421} base their security on the \gls{lwe} problem introduced by Regev \cite{regev:2009}; this thesis relies on its ring variant, \gls{rlwe} \cite{cryptoeprint:2012/230}.

An instantiation is fixed by a small set of parameters. The working ring $R_Q$ of \autoref{def:ring} contributes the ring dimension $N$ and the \emph{ciphertext modulus} $Q$. The \emph{plaintext modulus} $t \ll Q$ defines the plaintext space $R_t$, and the \emph{message scaling factor} $\Delta$ determines where inside a ciphertext the message is placed; its value is scheme-specific. Noise is drawn from an \emph{error distribution} $\chi$ over $R_Q$, a discrete Gaussian of small standard deviation $\sigma$. All parameters are chosen so that breaking the scheme requires on the order of $2^\lambda$ operations, where $\lambda$ denotes the \emph{security parameter}.

\begin{definition}[RLWE Distribution]
    \label{def:rlwe-dist}
    % TODO: §2.4 (gadgets/bv.tex, abstract gadget decomposition) calls a single ring element $v \in R_Q$ an "RLWE sample"; here a sample is a pair in $R_Q^2$ --- align terminology.
    For a fixed secret $s \in R_Q$, the \gls{rlwe} distribution $A_{s,\chi}$ over $R_Q^2$ is sampled by drawing $a \overset{\$}{\leftarrow} R_Q$ uniformly at random and $\epsilon \overset{\$}{\leftarrow} \chi$, and outputting the pair $(a \cdot s + \epsilon,\; a)$.
\end{definition}

\begin{definition}[Decision RLWE Problem]
    \label{def:rlwe-problem}
    Given arbitrarily many independent samples from $R_Q^2$, the decision \gls{rlwe} problem is to distinguish whether they were drawn from $A_{s,\chi}$ for a fixed secret $s$ or from the uniform distribution on $R_Q^2$. The \gls{rlwe} assumption states that no probabilistic polynomial-time adversary solves this problem with advantage better than negligible in $\lambda$ \cite{regev:2009, cryptoeprint:2012/230}.
\end{definition}

\begin{definition}[RLWE Ciphertext]
    \label{def:rlwe}
    An \gls{rlwe} ciphertext encrypting a plaintext $m \in R_t$ under a persistent ternary secret key $s \overset{\$}{\leftarrow} \{-1, 0, 1\}^N$ is a pair $x = (b, a) \in R_Q^2$, obtained by drawing a fresh sample from $A_{s,\chi}$ and adding the scaled message $\Delta m$ to its first component. The result consists of the \emph{masked element} $b$ and the \emph{public element} $a$.
    \begin{equation}
        \label{eq:rlwe}
        x = (b, a) = (a \cdot s + \Delta m + \epsilon,\; a)
    \end{equation}
    We write $\rlwe_s(m)$ for an encryption of $m$ under $s$, abbreviated $\rlwe(m)$ when the key is clear from context, and $\Call{Encrypt}{m} \mapsto \rlwe(m)$ for the procedure producing a fresh such ciphertext.
\end{definition}

Since fresh samples are computationally indistinguishable from uniform pairs by \autoref{def:rlwe-problem}, ciphertexts reveal nothing about their plaintexts; the formal IND-CPA reduction is given in \cite{cryptoeprint:2011/277}.

\begin{definition}[Decryption]
    \label{def:decrypt-generic}
    % TODO: 04-Results/noise.tex recalls this definition with the bound $\max_i \epsilon_i < Q$; the correct condition is the centered $Q/2$ bound of \autoref{def:decrypt-bgv} --- align.
    The function $\Call{Decrypt}{\rlwe_s(m), s} \mapsto m$ recovers the plaintext from a ciphertext $(b, a)$ by computing $b - a \cdot s = \Delta m + \epsilon$ and removing the noise $\epsilon$. Recovery is correct only while $\|\epsilon\|$ remains below a scheme-dependent bound.
\end{definition}

The linear operations on \gls{rlwe} ciphertexts are \emph{free}, in the sense that applying them to the ciphertext components applies them to the underlying plaintexts without any additional machinery.

\begin{definition}[Addition]
    \label{def:rlwe-add}
    The sum of two ciphertexts $\rlwe_s(m_x) = (b_x, a_x)$ and $\rlwe_s(m_y) = (b_y, a_y)$ under the same key is the component-wise sum, satisfying $\rlwe_s(m_x) + \rlwe_s(m_y) = \rlwe_s(m_x + m_y)$.
\end{definition}

Correctness follows by expanding the masked elements.
\begin{equation}
    b_x + b_y = (a_x + a_y) \cdot s + \Delta(m_x + m_y) + (\epsilon_x + \epsilon_y)
\end{equation}
The sum is a well-formed encryption of $m_x + m_y$ with public element $a_x + a_y$, but it carries the combined noise $\epsilon_x + \epsilon_y$.

\begin{definition}[Scalar Multiplication]
    \label{def:rlwe-scalar}
    % TODO: §2.4 (products.tex) types RGSW messages as $\mu \in R$, a ring never defined; settle on $R_Q$ or define $R = \mathbb{Z}[X]/(X^N+1)$ once.
    The product of a ciphertext $\rlwe_s(m) = (b, a)$ with a scalar $c \in R_Q$ is the component-wise product $c \cdot (b, a) = (c \cdot b,\; c \cdot a)$, satisfying $c \cdot \rlwe_s(m) = \rlwe_s(c \cdot m)$ with noise $c \cdot \epsilon$.
\end{definition}

Message and noise are scaled together, so only multipliers with small coefficients keep the result decryptable. Multiplication of two \emph{ciphertexts}, by contrast, requires considerable additional machinery and amplifies the noise sharply; no construction in this thesis performs it, and it is not treated further.

\begin{definition}[Public Key]
    \label{def:pk}
    A public key is an encryption of zero, $pk = \rlwe_s(0) = (a \cdot s + \epsilon_{pk},\; a)$. It enables encryption without knowledge of $s$: with fresh $u \overset{\$}{\leftarrow} \{-1, 0, 1\}^N$ and $e_0, e_1 \overset{\$}{\leftarrow} \chi$, the public-key encryption of $m \in R_t$ is defined below.
    \begin{equation}
        \label{eq:pk-encrypt}
        \Call{Encrypt}{m, pk} = u \cdot pk + (\Delta m + e_0,\; e_1)
    \end{equation}
    By \autoref{def:rlwe-add} and \autoref{def:rlwe-scalar} the result is a well-formed ciphertext $\rlwe_s(m)$ with noise $u \cdot \epsilon_{pk} + e_0 - e_1 \cdot s$, a sum of products of small elements.
\end{definition}

Grounding security in \gls{rlwe} has a structural cost: every ciphertext must carry a random noise term, and every homomorphic operation enlarges it. A ciphertext decrypts correctly only while its noise remains beneath the bound of \autoref{def:decrypt-generic}; once that budget is spent, the plaintext is irrecoverable. The noise budget is therefore the resource that determines how much computation a ciphertext can withstand.

Bootstrapping, due to Gentry \cite{gentry:2009}, lifts the restriction entirely: the decryption function is itself evaluated homomorphically, yielding a fresh encryption of the same plaintext in which the accumulated noise is replaced by the comparatively small noise of the bootstrapping procedure. The operation is expensive, and the leveled constructions in this thesis avoid it, instead choosing parameters large enough that the entire computation completes within a single budget. Within a budget, the standard instrument for shedding accumulated noise is the modulus switch.

\begin{definition}[Modulus Switching]
    \label{def:modswitch}
    Modulus switching transforms a ciphertext $x \in R_Q^2$ under the modulus $Q$ into a ciphertext $x' = \lfloor x \cdot Q'/Q \rceil \in R_{Q'}^2$ under a smaller modulus $Q' < Q$ that decrypts to the same plaintext, by scaling each component by $Q'/Q$ and rounding to the nearest integer. The noise shrinks by the same factor, making the operation the primary noise-management tool of leveled schemes.
\end{definition}

\subsection{The BGV Scheme}
\label{sec:bgv}
The implementation in this thesis is built on the BGV scheme \cite{cryptoeprint:2011/277}. %, chosen for its mature residue-based instantiation in OpenFHE \cite{cryptoeprint:2022/915}. 
As such, we introduce the primary differences between this and the general \gls{rlwe} setting, as far as we are concerned. %% TODO: Rewrite

BGV fixes the message scale to $\Delta = 1$ and scales every noise term by the plaintext modulus instead, so that the masked element of \autoref{def:rlwe} reads $b = a \cdot s + m + t\,\epsilon_{\mathsf{rlwe}}$ for a base noise term $\epsilon_{\mathsf{rlwe}} \overset{\$}{\leftarrow} \chi$. % \footnote{Some noise-management strategies, such as OpenFHE's \texttt{FLEXIBLEAUTO}, vary $\Delta$ in practice.} 
Ciphertexts are nevertheless written in the general form of \autoref{def:rlwe} throughout this thesis, with the substitution $\Delta = 1$, $\epsilon = t\,\epsilon_{\mathsf{rlwe}}$ left implicit.

\begin{definition}[BGV Decryption]
    \label{def:decrypt-bgv}
    BGV decryption removes the mask with a centered reduction modulo $Q$, and the noise with a reduction modulo $t$.
    \begin{equation}
        \Call{Decrypt}{(b, a), s} = \big[\, [b - a \cdot s]_Q \,\big]_t
    \end{equation}
    The inner reduction recovers $m + t\,\epsilon_{\mathsf{rlwe}}$ exactly over the integers provided $\|m + t\,\epsilon_{\mathsf{rlwe}}\| \leq \lfloor Q/2 \rfloor$, and the outer reduction then strips the noise term and returns $m$. This is the concrete instance of the bound in \autoref{def:decrypt-generic}: decryption fails once the noise reaches $Q/2$.
\end{definition}

\begin{definition}[\gls{simd} Packing]
    \label{def:simd}
    If the plaintext modulus is a prime satisfying $t \equiv 1 \pmod{2N}$, a plaintext may be reinterpreted as a vector of $N$ independent \emph{slots} $[c_i]_t$ on which homomorphic operations act element-wise \cite{cryptoeprint:2011/133}. For example, the slot-wise product gives $(1, 2) \cdot (2, 3) = (2, 6)$, in contrast to the polynomial product $(1 + 2X)(2 + 3X) = 2 + 7X + 6X^2$.
\end{definition}

\subsection{Multi-Party FHE}
\label{sec:mpfhe}
% TODO: bib: add a canonical threshold-FHE citation (e.g. Asharov et al., EUROCRYPT 2012); currently only the OpenFHE implementation is cited.
When mutually distrusting parties compute on shared data, no single party can be allowed to hold the secret key. Threshold variants of \gls{fhe} remove this single point of trust by distributing both key generation and decryption; this thesis uses the additive $n$-out-of-$n$ variant implemented in OpenFHE \cite{cryptoeprint:2022/915}.

During a one-time setup, each party $i$ samples its own secret share $s_i$ and contributes a corresponding public share, formed against a common public element $a$. By \autoref{def:rlwe-add} the shares aggregate into a single joint public key (\autoref{def:pk}) for the joint secret $s = \sum_i s_i$, under which every party encrypts. The joint secret itself is never reconstructed.

Decryption proceeds in two rounds. Each party first releases a \emph{partial decryption} $d_i = a \cdot s_i + \epsilon^{\mathsf{fl}}_i$ of the ciphertext $(b, a)$, computed from its own share and a fresh \emph{flooding} term. The flooding noise $\epsilon^{\mathsf{fl}}_i$ is drawn from a distribution much wider than $\chi$, so that $d_i$ statistically hides $s_i$. Because decryption is linear in the secret, the partial decryptions then combine into $b - \sum_i d_i = \Delta m + \epsilon - \sum_i \epsilon^{\mathsf{fl}}_i$ without any party learning the joint secret; correctness requires only that the flooding noise, like any other, stays within the noise budget of \autoref{def:decrypt-generic}. %% TODO: Cite
